% =============================================================
% Terimler Sözlüğü — bölümler ilerledikçe doldurulacak
% =============================================================

\textbf{Taban (Radix)} --- Pozisyonel bir sayı sisteminde kullanılan basamak sayısı; her
basamağın ağırlığı bu sayının kuvvetleriyle belirlenir (ör. ikili için $2$, onaltılık için
$16$).

\textbf{Ardışık Bölme/Çarpma Yöntemi} --- Ondalık bir sayıyı başka bir tabana çevirmek için
kullanılan yöntem; tam sayı kısmı ardışık bölme (kalanlar aşağıdan yukarı okunur), kesirli
kısım ardışık çarpma (taşan basamaklar yukarıdan aşağı okunur) ile çevrilir.

\textbf{Tümleyen (Complement)} --- Bir sayının çıkarma işlemini toplamaya çevirmek için
kullanılan dönüştürülmüş hâli; $(r{-}1)$'in tümleyeni (her basamağı $r{-}1$'den çıkararak) ve
$r$'nin tümleyeni (ona $1$ eklenerek) olmak üzere iki türü vardır.

\textbf{Uçtan-Dolaşan Elde (End-Around Carry)} --- $(r{-}1)$'in tümleyeniyle çıkarma
yapıldığında oluşan taşma bitinin atılmayıp sonucun en düşük anlamlı basamağına eklenmesi
kuralı.

\textbf{İşaret Biti} --- İşaretli bir ikili sayıda en soldaki bit; $0$ pozitif, $1$ negatif
anlamına gelir.

\textbf{Taşma (Overflow)} --- $n$ bitlik iki sayının toplamının $n$ bite sığmaması durumu;
donanımda sabit bit genişliği nedeniyle sonucun bozulmasına yol açar.

\textbf{Normalizasyon} --- Bir kayan noktalı sayıyı $1.xxxx\times2^k$ biçiminde, baştaki
basamak tam olarak $1$ olacak şekilde yazma işlemi; ikilide baştaki $1$ her zaman var olduğu
için saklanmaz (örtük bit).

\textbf{Bias} --- IEEE 754'te üs alanına eklenen sabit kaydırma değeri (tek duyarlıkta
$127$); negatif üsleri de işaretsiz tam sayı olarak karşılaştırılabilir kılar.

\textbf{Bool Değişkeni} --- Yalnızca $0$ veya $1$ değerini alabilen değişken.

\textbf{Mantık Kapısı} --- Bir veya birden çok ikili girişten tek bir ikili çıkış üreten
temel devre elemanı.
